\medskip








\begin{enumerate}
\item La solution $f$ de l'équation différentielle $y'=-3y + 7$ telle que $f(0) = 1$ est la fonction définie sur $\R$ par :

\smallskip

\begin{center}
\renewcommand\arraystretch{1.9}
\begin{tabularx}{\linewidth}{X X}
\textbf{A. } $f(x) = \e^{-3x}$&\textbf{B. } $f(x) = - \dfrac43 \e^{-3x} + \dfrac73$\\
\textbf{C. } $f(x) = \e^{-3x} + \dfrac73$&\textbf{D.~} $f(x) = - \dfrac{10}{3} \e^{-3x} - \dfrac73$
\end{tabularx}
\end{center}

L'équation différentielle $y' = - 3y$ a pour solutions les fonctions $x \longmapsto f(x) = K \e^{- 3x}$, avec $K \in \R$.

La fonction $x \longmapsto \alpha$, avec $\alpha \in \R$ est solution de l'équation différentielle $y' = - 3y + 7$ si et seulement si $y' = 0 = - 3\alpha + 7 \iff 3\alpha = 7 \iff \alpha = \dfrac73$.

On sait qu'alors les solutions de l'équation différentielle $y' = - 3y + 7$ sont les fonctions $x \longmapsto K \e^{- 3x} + \dfrac73$.

En particulier la fonction $f$ solution telle que $f(0) = 1 \iff K + \dfrac73 = 1 \iff K = - \dfrac43$.

La seule solution est donc la fonction définie par $f(x) = - \dfrac43\e^{-3x} + \dfrac73$ : réponse \textbf{B.}
\item La courbe d'une fonction $f$ définie sur $[0~;~+\infty[$ est donnée ci-dessous.
	
\begin{center}
\psset{unit=1.4cm}
\begin{pspicture}(-0.5,-0.25)(6.25,4.25)
\pscustom[fillstyle=solid,fillcolor=lightgray]{\psline(1,0)(1,3)\psplot[plotpoints=2000,linewidth=1.25pt,linecolor=red]{1}{5}{0.00196  x 5 exp mul 0.05764 x  4 exp mul sub 0.62643 x  3 exp mul add 3.00524 x dup mul mul sub 5.43449  x mul add}\psline(5,0.4)(5,0)}
\psgrid[gridlabels=0pt,subgriddiv=1]
\psaxes[linewidth=1.25pt,labelFontSize=\scriptstyle]{->}(0,0)(0,0)(6.25,4.25)
\psplot[plotpoints=2000,linewidth=1.25pt,linecolor=red]{0}{6.25}{0.00196  x 5 exp mul 0.05764 x  4 exp mul sub 0.62643 x  3 exp mul add 3.00524 x dup mul mul sub 5.43449  x mul add}
\pspolygon[fillstyle=hlines](1,0)(1,3)(2,3)(2,2)(3,2)(3,0)
\pspolygon[linecolor=blue,linewidth=1.25pt](1,0)(1,4)(2,4)(2,3)(3,3)(3,2)(4,2)(4,1)(5,1)(5,0)
\end{pspicture}
\end{center}









Le dessin est clair : la fonction est positive sur l'intervalle[1~;~5] ; l'intégrale est (en unités d'aire) la mesure de la surface limitée par la représentation graphique de $f$, l'axe des abscisses et les droites d'équations $x = 1$ et $x = 5$.

La surface grise contient les 5 carreaux hachurés et est inscrite dans le polygone de 10 unités en bleu. Réponse \textbf{C.}
\item 









On sait que si $g'$ est la dérivée de $g$, alors $g$ est une primitive de la fonction $g'$, donc :

$\displaystyle\int_0^2 g'(x)\:\text{d}x = \left[x^2 \ln \left(x^2 + 4\right)\right]_0^2 = 2^2\ln \left(2^2 + 2 \right) = 4\ln 4 + 4)  = 4\ln 8$ ou $4\ln 2^3 = 3 \times 4 \ln 2 = 12\ln 2 \approx 8,31$ soit 8,3 au dixième près. Réponse \textbf{B.}
	\smallskip
	
\item 
	









Le nombre de groupes de 5 élèves parmi les 31 est $\displaystyle\binom{31}{5}$. Réponse : \textbf{D.}
	\item 
	








	







Elle choisit 3 élèves parmi les 20 faisant SES : elle a $\displaystyle\binom{20}{3}$ possibilités ; ensuite dans chacun de ces cas elle choisit 2 élèves parmi les $31 - 20 = 11$ élèves qui ne font pas SES, ce qui fait $\displaystyle\binom{20}{3} \times \binom{11}{2}$ possibilités. Réponse \textbf{A.}
\end{enumerate}

\bigskip
